\documentclass[11pt, a4paper]{article}
% === Libertine (Linux Libertine) ===
\usepackage[top=0.6in, bottom=0.6in, left=0.8in, right=0.8in]{geometry}
\usepackage[T1]{fontenc}
\usepackage{libertine}
\usepackage{microtype}
\usepackage{fontawesome5}
\usepackage{setspace}
\setlength{\parindent}{0pt}
\setstretch{1.3}
\begin{document}

{\fontsize{24pt}{30pt}\selectfont\bfseries Linsheng He}

\vspace{6pt}
{\LARGE Libertine \textbullet\ Linux Libertine}

\vspace{12pt}
{\footnotesize 识别特征：}\\[2pt]
{\large \textbf{Bold} \textit{Italic} \textsc{Small Caps} Rg Regular}\\[4pt]
{\large Policy design, behavioral public management, government communication.}\\[4pt]
{\large\itshape Public Administration, Review of Policy Research.}

\vspace{16pt}
{\footnotesize\bfseries 字号对比：}\\[2pt]
{\tiny Linsheng} {\scriptsize Linsheng} {\footnotesize Linsheng} {\small Linsheng}
Linsheng {\large Linsheng} {\Large Linsheng}

\vspace{12pt}
\begin{itemize}
    \item f --- 弧线尾端有钩形装饰，这是 Libertine 最明显特征
    \item t --- 底部向右弯曲
    \item g --- 底环不完全闭合，有缺口
    \item Libertine 特点：\textbf{温暖}\ 且 \textbf{有个性}，不像传统学术体那么冷
\end{itemize}

\end{document}
